\section{The Witness of the Fathers Before and After Dionysius}
\label{sec:patristic-witness}

Around the year~350, instructing newly baptized Christians inside the Church of the Holy Sepulchre in Jerusalem, St.~Cyril unfolded the meaning of the Eucharistic Preface they had just prayed for the first time:
\begin{quotation}
After this, we make mention of heaven, and earth, and sea; of sun and moon; of stars and all the creation, rational and irrational, visible and invisible; of Angels, Archangels, Virtues, Dominions, Principalities, Powers, Thrones; of the Cherubim with many faces: in effect repeating that call of David's \emph{Magnify the Lord with me}. We make mention also of the Seraphim, whom Esaias in the Holy Spirit saw standing around the throne of God... crying \emph{Holy, Holy, Holy, is the Lord of Sabaoth}. For the reason of our reciting this confession of God, delivered down to us from the Seraphim, is this, that so we may be partakers with the hosts of the world above in their Hymn of praise. (St.~Cyril of Jerusalem, \emph{Mystagogical Lecture} 5.6)
\end{quotation}
A century and a half before Pseudo-Dionysius composed \emph{The Celestial Hierarchy}, the worshipping Church already named all nine scriptural choirs in a single breath at the altar. Patristic angelology did not begin as abstract speculation; it grew out of the Eucharist, out of martyrdom and combat with pagan idolatry, out of anti-Gnostic and anti-Arian controversy, and out of close meditation upon Genesis, the Psalms, the Prophets, and St.~Paul.

\subsection{Apostolic and Pre-Nicene Fathers}
\label{subsec:pre-nicene-fathers}

The earliest Christian voice after the Apostles already joins the earthly assembly to the angels' unceasing worship. Writing from Rome to heal faction in Corinth near the close of the first century, \textbf{St.~Clement of Rome} weaves Daniel~7:10 and Isaias~6:3 together into the oldest surviving extra-canonical testimony to the liturgical \emph{Sanctus}, holding up angelic harmony as the rule for the Church on earth:
\begin{quotation}
Let us consider the whole multitude of His angels, how they stand ever ready to minister to His will. For the Scripture saith, ``Ten thousand times ten thousand stood around Him, and thousands of thousands ministered unto Him, and cried, Holy, holy, holy, [is] the Lord of Sabaoth; the whole creation is full of His glory.'' And let us therefore, conscientiously gathering together in harmony, cry to Him earnestly, as with one mouth. (St.~Clement of Rome, \emph{1~Clement} 34.5--7)
\end{quotation}

In the second century, \textbf{St.~Justin Martyr} brought Christian belief about the angels into direct conversation with both Roman paganism and Jewish Scripture. Refuting the charge of atheism in his \emph{First Apology}~(6), Justin wrote a sentence whose syntax has been closely examined ever since:
\begin{quotation}
But both Him, and the Son (who came forth from Him and taught us these things, and the host of the other good angels who follow and are made like to Him), and the prophetic Spirit, we worship and adore, knowing them in reason and truth. (St.~Justin Martyr, \emph{First Apology}~6)
\end{quotation}
Some readers construe ``the host of the other good angels'' as joined with the Father, Son, and Spirit in a broad sense of reverential honor---never infringing the strict monotheism Justin defends everywhere else---while others take the angels as an object of what the Son taught us. Either way, Justin places the holy angels firmly on the side of those who follow and are assimilated to the Son. In his \emph{Second Apology}~(5), as we saw earlier, he received the Enochic Watchers reading of Genesis~6 to argue that fallen angels and their demonic offspring lay behind the myths of the pagan gods. Most influential for early Christology was his \emph{Dialogue with Trypho}~(56--61, 128), where Justin argued that the divine Person appearing to Abraham at Mamre, judging Sodom with two attendant angels, and speaking from the burning bush was the pre-incarnate Logos Himself: ``He who appeared to Abraham under the oak in Mamre is God, sent with the two angels in His company to judge Sodom by Another who remains ever in the supercelestial places'' (\emph{Dialogue}~56). Against modalists who reduced the Son to a fleeting projection, Justin insisted that Christ is called ``Angel'' because He announces the Father's counsel while remaining personally distinct (\emph{Dialogue}~128)---a Christological reading of Old Testament theophanies that post-Nicene Fathers later refined once Arianism twisted the title ``Angel'' to imply creaturehood.

\textbf{St.~Irenaeus of Lyons} forged the enduring anti-Gnostic foundation of Catholic angelology around A.D.~180. Valentinian Gnostics attributed the creation of the visible cosmos and of human bodies to ignorant or rebellious angels, cleaving the supreme Father from the world of matter. Irenaeus dismantled that dualism at the root in \emph{Against Heresies}. Naming ``Principalities, Powers, Angels, Archangels, Dominations, Virtues''---an early, unranked gathering of Pauline titles---he declared: ``we prove from the authoritative Scriptures that all the things which have been mentioned, visible and invisible, have been made by one God'' (\emph{Against Heresies} II.30.6). God never delegated creation to angelic demiurges, because He acts immediately by His own ``two hands,'' the Son and the Holy Spirit:
\begin{quotation}
the Father being in no want of angels, in order that He might call the creation into being, and form man... For His offspring and His similitude do minister to Him in every respect; that is, the Son and the Holy Spirit, the Word and Wisdom; whom all the angels serve, and to whom they are subject. (St.~Irenaeus, \emph{Against Heresies} IV.7.4; cf.~IV.20.1)
\end{quotation}
With equal vigor, Irenaeus anchored demonic evil strictly in rational freedom rather than in a corrupted nature: ``And in man, as well as in angels, He has placed the power of choice (for angels are rational beings), so that those who had yielded obedience might justly possess what is good, given indeed by God, but preserved by themselves'' (\emph{Against Heresies} IV.37.1). God made neither angels nor men wicked by nature; the devil created nothing whatsoever, ``since indeed he is himself a creature of God, like the other angels,'' and eternal fire is prepared for him as ``the ringleader of the apostasy'' and for those who revolted with him (IV.40.1, IV.41.1; cf.~III.23.3).

In third-century Alexandria, \textbf{Origen} opened questions that shaped both Greek and Latin theology for a millennium, even where the Church later rejected his boldest hypotheses. In his \emph{Commentary on Matthew} (XIII.26--29), meditating on Christ's words concerning the angels of the ``little ones'' (Matt~18:10), Origen posed the exact question St.~Thomas Aquinas would later settle in \ST{I, q.~113, a.~5}: when does a guardian angel assume charge of a human soul---at baptism, through the laver of regeneration, or from birth according to God's foreknowledge and predestination? Origen set forth the orthodox view that Christ, having redeemed us by His blood, entrusts the baptized believer to a holy angel whose purity fits him to behold the Father's face. Yet beside that view he floated a speculative third proposal: perhaps an evil angel assigned at birth might gradually come to believe alongside the human soul and advance until he too beheld the face of the Father. The Church definitively rejected that speculation and the wider Origenist cosmology of pre-existent minds falling into angelic, human, and demonic ranks (Constantinople~543 and~553), while retaining Origen's penetrating inquiry into guardianship and spiritual illumination.

\subsection{The Fourth-Century East}
\label{subsec:fourth-century-east}

During the Arian and Macedonian controversies of the fourth century, the Eastern Fathers defined the boundary between the uncreated Trinity and the created angels with lasting sharpness, while deepening the Church's theology of angelic nature, grace, and discernment.

\textbf{St.~Athanasius of Alexandria} confronted Arian attempts to place the eternal Son at the summit of the angelic ranks. Commenting on Hebrews~1:4--6 in his first discourse \emph{Against the Arians}~(I.61--62), he drew an unbridgeable line between every ministering spirit and the begotten Word:
\begin{quotation}
And the Angels indeed minister ascending and descending; but concerning the Son he saith, ``And let all the Angels of God worship Him.'' And when Angels minister, they say, ``I am sent unto thee,'' and, ``The Lord has commanded;'' but the Son... says, as being Word and Image, ``I am in the Father, and the Father in Me.'' (St.~Athanasius, \emph{Against the Arians} I.61--62)
\end{quotation}
Alongside that Christological clarity, Athanasius bequeathed to Christian ascetical tradition its classic rule for discerning angelic from demonic visitations. In the \emph{Life of Antony}~(\S 35), he records the desert father's teaching:
\begin{quotation}
For the presence either of the good or evil by the help of God can easily be distinguished. The vision of the holy ones is not fraught with distraction... But it comes so quietly and gently that immediately joy, gladness and courage arise in the soul... But if, being men, some fear the vision of the good, those who appear immediately take fear away; as Gabriel did in the case of Zacharias, and as the angel did who appeared to the women at the holy [sepulchre]. (St.~Athanasius, \emph{Life of Antony}~\S 35)
\end{quotation}
Demonic apparitions, by contrast, arrive with crash and tumult and leave behind lasting terror, confusion, and dejection---a criterion St.~Thomas Aquinas will cite directly when analyzing Gabriel's greeting to Mary (\ST{III, q.~30, a.~3, ad~3}).

In Cappadocia, \textbf{St.~Basil the Great} framed the classical Eastern synthesis of angelic creation, sanctification, and guardianship. Preaching on Genesis~1:1 in the \emph{Hexaemeron}~(I.5), Basil argued that an invisible intellectual realm existed even before the visible universe and its time began:
\begin{quotation}
It appears, indeed, that even before this world an order of things existed... outstripping the limits of time, eternal and infinite. The Creator and Demiurge of the universe perfected His works in it, spiritual light for the happiness of all who love the Lord, intellectual and invisible natures, all the orderly arrangement of pure intelligences who are beyond the reach of our mind and of whom we cannot even discover the names. They fill the essence of this invisible world, as Paul teaches us. ``For by him were all things created... whether they be thrones or dominions or principalities or powers'' or virtues or hosts of angels or the dignities of archangels. (St.~Basil, \emph{Hexaemeron} I.5)
\end{quotation}
Basil's reverent reticence---remarking that many celestial orders bear names beyond human discovery---reappears in his treatise \emph{On the Holy Spirit}~(XVI.38), where he counters Pneumatomachians who degraded the Holy Spirit to a ministering spirit. Angels are not holy by their unaided essence; their holiness is the gift of the Spirit, and their whole existence bears the single imprint of the Trinity:
\begin{quotation}
The pure, intelligent, and supermundane powers are and are styled holy, because they have their holiness of the grace given by the Holy Spirit... And in the creation bethink thee first, I pray thee, of the original cause of all things that are made, the Father; of the creative cause, the Son; of the perfecting cause, the Spirit; so that the ministering spirits subsist by the will of the Father, are brought into being by the operation of the Son, and perfected by the presence of the Spirit. Moreover, the perfection of angels is sanctification and continuance in it. (St.~Basil, \emph{On the Holy Spirit} XVI.38)
\end{quotation}
Turning to the ministry of angels among men, Basil meditates in his \emph{Homily on Psalm~33} on the angel who encamps around those who fear God, and sets down his most celebrated declaration on personal guardianship in \emph{Against Eunomius}~(III.1): though all angels share one spiritual nature, some preside over nations while others attend individual believers, so that ``no one will contradict that beside every one of the faithful stands an angel, as a kind of tutor and shepherd guiding his life'' (\latin{velut paedagogus quidam et pastor vitam dirigens}), citing the Lord's words in Matthew~18:10.

Basil's friend \textbf{St.~Gregory of Nazianzus} gave Greek angelology its most luminous philosophical language. In his Second Theological Oration (\emph{Oration} 28.31), expounding Psalm~103:4, Gregory notes that the angel is called ``Spirit, as being a creature of the intellectual sphere; Fire, as being of a purifying nature,'' and adds with careful modesty that ``relatively to us at least, we must reckon the Angelic Nature incorporeal, or at any rate as nearly so as possible.'' Contemplating their multitude until human reason grows dizzy, he gathers both scriptural ranks and contemplative epithets---``Angels and Archangels, Thrones, Dominions, Princedoms, Powers, Splendours, Ascents, Intelligent Powers or Intelligencies''---and coins the phrase that echoes down through Dionysius and Damascene: ever circling the First Cause, they are ``so conformed to beauty and moulded that they become secondary Lights, and can enlighten others by the overflowings and largesses of the First Light.'' In his Nativity sermon (\emph{Oration} 38.9--10, repeated verbatim in \emph{Oration} 45.5 on Easter), Gregory traces the Trinitarian creation of these secondary splendors prior to the material world and faces the sobering fact of the fall:
\begin{quotation}
He first conceived the Heavenly and Angelic Powers. And this conception was a work fulfilled by His Word, and perfected by His Spirit. And so the secondary Splendours came into being, as the Ministers of the Primary Splendour... I should like to say that they were incapable of movement in the direction of evil, and susceptible only of the movement of good... but I am obliged to stop short of saying that, and to conceive and speak of them only as difficult to move because of him, who for his splendour was called Lucifer, but became and is called Darkness through his pride; and the apostate hosts who are subject to him, creators of evil by their revolt against good and our inciters. (St.~Gregory of Nazianzus, \emph{Oration} 38.9--10)
\end{quotation}

\textbf{St.~Gregory of Nyssa} extended Cappadocian inquiry to the puzzle of angelic multiplicity. Answering how humanity might have multiplied in paradise without marriage in \emph{On the Making of Man}~(17.2--3), he pointed to the angels, asking ``what is the mode of the angelic existence, how they exist in countless myriads, being one essence, and at the same time numerically many''---a question whose metaphysical tension between unity of essence and numerical plurality would later divide medieval Thomists and Franciscans.

At Antioch and Constantinople, \textbf{St.~John Chrysostom} joined rigorous anti-Arian exegesis to keen pastoral preaching. Expounding Matthew~18:10 (\emph{Homily~59 on Matthew}~4), he appealed to 1~Corinthians~11:10 and Deuteronomy~32:8 LXX to conclude: ``Hence it is evident, that the saints have angels, or even all men... But here He is discoursing not of angels only, but rather of angels that are greater than others. But when He saith, `The face of my Father,' He means nothing else than their fuller confidence, and their great honor''---a remark St.~Thomas Aquinas will carefully reconcile with the Dionysian assignment of personal guardianship to the lowest choir (\ST{I, q.~113, a.~3, ad~1}). Preaching on Ephesians~1:21 (\emph{Homily~3 on Ephesians}), Chrysostom celebrates Christ's exalted humanity placed above every heavenly rank, adding with Basilian apophatic restraint, ``Therefore powers there are whose names are to us unintelligible, and unknown.'' On Colossians~1:16 (\emph{Homily~3 on Colossians}), he shuts the door against any subordination of the Spirit or independence of angelic creatures: ``But the Spirit is not amongst the `powers.'... Not only did He Himself bring them out of nothing into being, but Himself sustains them now, so that were they dissevered from His Providence, they were at once undone and destroyed.''

\subsection{The Western Patristic Synthesis}
\label{subsec:western-fathers}

In the Latin West, angelology unfolded through biblical commentary, anti-Arian polemic, and above all the immense architecture of St.~Augustine.

\textbf{St.~Hilary of Poitiers}, commenting on Psalm~124, set angelic guardianship firmly in its proper place beneath God's own presence. Recalling how God offered Moses an angel to go before Israel (Ex~32:34) and how Moses answered that he would not go up unless the Lord Himself walked with them (Ex~33:15), Hilary summed up the relation in seven Latin words: \latin{Bonum quidem praesidium angeli, sed melius Domini}, ``Good indeed is the protection of an angel, but better that of the Lord.''

\textbf{St.~Ambrose of Milan} marshaled the heavenly choirs in defense of both Nicene Christology and the full divinity of the Holy Spirit. Challenging Arian judges in \emph{On the Faith}~(V.18.225), he asked: ``Who shall decide about the Son of God? Thrones, dominions, angels, powers? But archangels give attendance and serve Him, cherubim and seraphim minister to Him and praise Him.'' In \emph{On the Holy Spirit}~(I.16.178), he pictured the Spirit proceeding from the Fount of Life to ``flow more abundantly among those celestial Thrones, Dominions and Powers, Angels and Archangels, rushing in the full course of the seven virtues of the Spirit,'' and pressed the anti-Macedonian conclusion in Book~III~(16.142): ``For it is said of the Son of God: `Let all the Angels of God worship Him;' we do not read, Let the Spirit worship Him. For how can He worship Who is not amongst servants and ministers?'' Throughout his treatises on the Spirit and on the Six Days (\emph{Hexameron}), Ambrose contrasted the uncircumscribed immensity of God with the circumscribed finitude of every created angel, while his \emph{Exposition of the Gospel of Luke}~(Book~II) meditated with Marian devotion on Gabriel's mission to Nazareth and the divine wisdom veiled in Mary's espousal.

\textbf{St.~Jerome} brought Hebrew scholarship and sharp polemical precision to the scriptural texts on the angels. Expounding Daniel~10:13 in his \emph{Commentary on Daniel}, Jerome joined the ``prince of the kingdom of the Persians'' directly to Deuteronomy~32:8 LXX, explaining that this prince appeared to be the good guardian angel entrusted with Persia, pleading for the province committed to him lest all the Jewish captives depart, until Michael---who presides over Israel---came to help; and Jerome defined the title ``chief princes'' concisely: \latin{Principes autem primos, archangelos intelligimus}, ``By chief princes we understand archangels.'' In his exegetical work on Isaias, Jerome guarded the two six-winged seraphim of Isaias~6:2 against an Origenist reading that had hazardously likened them to the Son and the Holy Spirit. And in his commentary on Matthew~18:10, Jerome uttered the sentence that St.~Thomas Aquinas places at the threshold of his question on guardian angels (\ST{I, q.~113, a.~2}): ``Great is the dignity of souls, for each one to have an angel deputed to guard it from its birth.''

\textbf{St.~Augustine of Hippo} gave the Latin Church the most comprehensive patristic theology of the angels before Aquinas, woven across \emph{The City of God} (Books~VIII--XII, XIV--XV), the \emph{Literal Commentary on Genesis}, and \emph{On the Trinity} (Book~III). Having refuted pagan demonology and Iamblichean theurgy in Books~VIII--X of \emph{The City of God}, Augustine turns in Books~XI and~XII to the origin, knowledge, beatitude, and fall of the angels.

First, confronting the silence of Genesis~1 regarding the explicit word ``angels,'' Augustine refuses to believe Moses omitted them, since God rested on the seventh day from all His works (\emph{City of God} XI.9). He identifies their creation on the first day under the name of ``heaven,'' or more fittingly under the name of ``light'' (\latin{Fiat lux}): ``if mention of them is made, it is implicitly under the name of `heaven,' when it is said, `In the beginning God created the heavens and the earth,' or perhaps rather under the name of `light.''' That reading shaped the Latin consensus that angels were created simultaneously with the visible cosmos at the beginning of time---a position distinct from the Eastern view of Basil, Gregory of Nazianzus, and John of Damascus, and later mirrored in the wording of Lateran~IV.

Second, Augustine read the refrain of Genesis~1---``and there was evening and morning''---as a figure of two modes of angelic cognition, \textbf{morning knowledge} and \textbf{evening knowledge} (\emph{City of God} XI.29; \emph{Literal Commentary on Genesis} IV.31.48). Because the holy angels behold creatures first in the wisdom of God as in the creative art by which they were made, they enjoy a brighter ``noonday'' knowledge in the Word than their ``twilight'' knowledge of creatures (and of themselves) in their own created nature---a contrast St.~Thomas Aquinas examines in full in Section~\ref{sec:st-knowledge}. In the \emph{Literal Commentary on Genesis} (IV.31.48), Augustine adds that when the angels know a created work first in the Word of God and then in its own proper nature (evening), that evening turns once more into morning (\latin{facta vespera fiebat mane}) when they refer the creature back to the praise of its Creator---whereas remaining fixed in creaturely self-regard becomes night.

Third, Augustine resolved the mystery of the angelic fall without yielding an inch to Manichaean dualism. Good and fallen angels share one nature made wholly good by God, differing solely ``from a difference in their wills and desires'' (\emph{City of God} XII.1). Asking what caused the evil will of the first apostate angel, Augustine answers that an evil will has no efficient cause, only a \emph{deficient} cause (\emph{City of God} XII.6, 9): ``they are miserable because they have forsaken Him who supremely is, and have turned to themselves who have no such essence. And this vice, what else is it called than pride?''---a purely spiritual sin demonstrating that the root of sin lies in the proud will, never in bodily flesh (\emph{City of God} XIV.3). In \emph{City of God} XII.9, Augustine holds it most fitting that God created the angels already endowed with good will---``as soon as ever they were created, they attached themselves to Him who created them, with the love He created in them.'' Yet in \emph{Literal Commentary on Genesis} (XI.16.21), glossing John~8:44, he maintains that the devil apostatized from the very dawn of creation (\latin{ab initio temporis}), never enjoying true blessedness, because genuine beatitude demands both uninterrupted fruition of God and unshakable certainty of eternal perseverance: ``that it be delivered from all dubiety, and know certainly that it shall eternally abide in the same enjoyment'' (\emph{City of God} XI.13).

Fourth, in \emph{On the Trinity} (III.4--9), Augustine lays down the metaphysical law governing angelic action upon the material world. God alone creates new natures out of nothing. Neither holy angels nor demons can be called creators when startling bodily changes occur at their agency---not even when Pharaoh's magicians produced frogs and serpents. Rather, all lower bodies are governed in due order by spiritual creatures under divine providence, and angels work visible wonders by swiftly gathering and applying the hidden natural ``seeds'' (\latin{rationes seminales}) that God Himself implanted within the corporeal elements of the world, always strictly within the bounds of divine permission.

Closing the fifth-century West, \textbf{Pope St.~Leo the Great} wrote no systematic catalogue of angelic choirs, yet his Nativity and Ascension sermons set angelology squarely around the mystery of the Word made flesh. ``The mystery of the Incarnation is a fitting theme for joy both to angels and to men'' (\emph{Sermon} XXI.2), because the Son's humility heals the breach between heaven and earth; and on Ascension Day (\emph{Sermon} LXXIV.1, quoted in full in Section~\ref{sec:christ-mary-angels}), Leo lifts the Church's gaze above every angelic host and rank of powers to behold our own humble human nature seated in Christ at the right hand of the Father.

\subsection{Post-Dionysian Patristic Culmination}
\label{subsec:post-dionysian-fathers}

Between the sixth and eighth centuries, two Doctors---one in Rome, one in Palestine---received the Dionysian corpus and gave Catholic angelology its mature patristic form in West and East.

In June~591, preaching in the Basilica of Saints John and Paul in Rome on the Gospel of the lost sheep and the lost drachma (Luke~15:1--10), \textbf{Pope St.~Gregory the Great} delivered his celebrated \emph{Homily~34 on the Gospels} (\S\S 3, 6--14), the single most influential Latin patristic exposition of the angelic orders. Why did the woman in the parable hold ten drachmas? Because ``there are nine orders of angels; yet that the number of the elect might be filled up, man was created as the tenth'' (\S 6). Combining Ephesians~1:21, Colossians~1:16, the prophetic books, and the nine precious stones of Ezechiel~28:13, Gregory enumerates the nine choirs in ascending order (\S 7): \latin{angelos, archangelos, virtutes, potestates, principatus, dominationes, thronos, cherubim, atque seraphim}---Angels, Archangels, Virtues, Powers, Principalities, Dominations, Thrones, Cherubim, and Seraphim. (Notably, in his earlier \emph{Moralia in Job} XXXII.23.48, commenting on the same nine stones of Ezechiel, Gregory had printed a different sequence: ``Angels, Archangels, Thrones, Dominations, Virtues, Princedoms, Powers, Cherubim, and Seraphim,'' placing Thrones beneath Dominations and Powers above Virtues and Principalities.)

In \S\S 8--10 of Homily~34, Gregory distills the lexical and functional theology of the angels into formulas memorized by every medieval master. First, he sharpens Augustine's distinction between being and mission:
\begin{quotation}
\latin{Sciendum quoque quod angelorum vocabulum, nomen est officii, non naturae. Nam sancti illi coelestis patriae spiritus semper quidem sunt spiritus, sed semper vocari angeli nequaquam possunt, quia tunc solum sunt angeli, cum per eos aliqua nuntiantur.} (St.~Gregory the Great, \emph{Homily~34 on the Gospels}~\S 8)
\end{quotation}
``We must also know that the word `angel' is a name of office, not of nature. For those holy spirits of the heavenly country are indeed always spirits, yet they can by no means always be called angels, because they are angels only when something is announced through them.'' Lesser tidings belong to Angels, supreme announcements to Archangels, whose very scriptural names match their missions: ``Michael means `Who is like God'; Gabriel, `Strength of God'; Raphael, `Medicine of God''' (\S 9). Turning to the seven higher orders (\S 10), Gregory defines each by its characteristic gift: \textbf{Virtues} (\latin{Virtutes}) are spirits through whom signs and miracles are most often worked; \textbf{Powers} (\latin{Potestates}) hold hostile forces subject to their sway so that demons cannot tempt human hearts as fiercely as they wish; \textbf{Principalities} (\latin{Principatus}) preside even over good angelic spirits, ordering what each is to carry out; \textbf{Dominations} (\latin{Dominationes}) excel in wondrous authority to which the rest are subject in obedience; \textbf{Thrones} (\latin{Throni}) are those in whom Almighty God ever presides to exercise judgment; \textbf{Cherubim} signifies ``fullness of knowledge'' (\latin{plenitudo scientiae}) drawn from nearness to God; and \textbf{Seraphim} means ``burning or kindling'' (\latin{ardentes vel incendentes}), spirits so intimately joined to God that no other spirits stand between them and Him, burning the more ardently the closer they behold Him. In \S 11 Gregory maps all nine ranks onto nine vocations among the elect in the Church, and in \S 14 he crowns the social vision of heaven with the communion of gifts: in that supreme city, what each order holds in part, all possess in common, yet each choir takes its proper name from the gift it has received most abundantly.

What makes Homily~34 decisive in the history of Western angelology is \S\S 12--13, where Gregory becomes the first Latin Father to cite Pseudo-Dionysius explicitly by name: \latin{Fertur vero Dionysius Areopagita, antiquus videlicet et venerabilis Pater}, ``Dionysius the Areopagite, an ancient and venerable Father, is said to declare that from the lesser ranks of angels messengers are sent forth visibly or invisibly to fulfill external ministry'' (\S 12), since it is one thing to minister and another to assist before the throne. Yet Gregory immediately addresses how an angel sent on an earthly errand can still ``always see the face of my Father'' (Matt~18:10), answering with one of the finest sentences in Christian theology:
\begin{quotation}
\latin{Et mittuntur igitur, et assistunt, quia etsi circumscriptus est angelicus spiritus, summus tamen spiritus ipse, qui Deus est, circumscriptus non est. Angeli itaque et missi, et ante ipsum sunt, quia quolibet missi veniant, intra ipsum currunt.} (St.~Gregory the Great, \emph{Homily~34 on the Gospels}~\S 13)
\end{quotation}
``Therefore they are both sent and stand before Him; for although the angelic spirit is circumscribed, yet the supreme Spirit Himself, who is God, is uncircumscribed. And so the angels are at once sent and in His presence, because wherever they go when sent, they run within Him.''

A century and a half later, in the monastery of Mar Saba near Jerusalem, \textbf{St.~John of Damascus} gathered the Greek patristic tradition into a crystalline synthesis in his \emph{Exposition of the Orthodox Faith} (Book~II, chs.~2--4). His formal definition became standard in both East and West:
\begin{quotation}
An angel, then, is an intelligent essence, in perpetual motion, with free-will, incorporeal, ministering to God, having obtained by grace an immortal nature... For all that is compared with God Who alone is incomparable, we find to be dense and material. For in reality only the Deity is immaterial and incorporeal. (St.~John of Damascus, \emph{Exposition of the Orthodox Faith} II.3)
\end{quotation}
This relative incorporeality---echoing Gregory of Nazianzus---guards God's unique transcendence without attributing bodily extension to angels. Indeed, John immediately specifies how angels occupy place: ``Seeing that they are minds they are in mental places, and are not circumscribed after the fashion of a body. For they have not a bodily form by nature, nor are they extended in three dimensions. But to whatever post they may be assigned, there they are present after the manner of a mind and energise, and cannot be present and energise in various places at the same time'' (II.3). From that same incorporeality he derives the irrevocability of the demonic fall: ``It is not susceptible of repentance because it is incorporeal. For it is owing to the weakness of his body that man comes to have repentance'' (II.3).

John adopts the Dionysian ninefold hierarchy explicitly by name---``that most holy, and sacred, and gifted theologian, Dionysius the Areopagite''---reproducing all three triads in exact descending order: Seraphim, Cherubim, Thrones; Dominions, Powers (\gk{Dynameis}), Authorities (\gk{Exousiai}); Rulers (\gk{Archai}), Archangels, Angels (II.3). On the moment of creation, he sides explicitly with Gregory of Nazianzus against Augustine's simultaneous Day-One reading: ``For myself, I am in harmony with the theologian. For it was fitting that the mental essence should be the first created, and then that which [is perceptible by sense]'' (II.3). Finally, turning to the fallen angels (II.4), he teaches that ``He who from among these angelic powers was set over the earthly realm... was not made wicked in nature but was good, and made for good ends,'' fell solely by turning from good to privation---``just as darkness also is absence of light''---and exercises no power over any creature ``except what God in His dispensation hath conceded to them.'' With Gregory in the West and John of Damascus in the East, the patristic synthesis stood ready to meet the systematic rigor of the medieval schools---and at the center of both stood the enigmatic treatise of Pseudo-Dionysius.
